\section{Conclusion}\label{sec:conclusion}

This work established and validated a complete, open-source computational ecosystem in \texttt{ngspice-47} and Python for behavioral modeling, physical characterization, hardware security evaluation, and visual prototyping of digital circuits based on two-dimensional BP SB-FET devices. Translating the 4D experimental surfaces with bilateral source referencing and gate-specific saturation curves (\texttt{bp\_ari\_ngspice.lib} \texttt{v2}) eliminated the proprietary dependency on Synopsys HSPICE/Verilog-A with mean error below $0.0621\%$ against the source tables, achieved $R^2 \ge 0.99977$ across \texttt{N\_transfer}, \texttt{P\_transfer}, and \texttt{P\_output}, and reduced \textit{p-FET} output RMSE by $196\times$ by explicitly capturing the measured $V_{g2} = -1.5\ \text{V}$ regime.

At the circuit and security levels, electrostatic $V_{fb}$ compensation preserved differential voltage gain above unity ($|A_v| = 2.71$) and $89.75\%$ output swing down to $V_{DD} = 0.2\ \text{V}$ in the complementary inverter. The 4-transistor \texttt{NAND}/\texttt{NOR} and \texttt{XOR}/\texttt{XNOR} polymorphic primitives attained $100\%$ dynamic truth-table accuracy with $0\%$ transistor overhead and $100.0\%$ yield across $250$ Monte Carlo supply-rail noise samples. In the 6T SRAM cell, Mechanism B was validated under fixed power rails: tuning $V_{g2,\mathrm{PU}} = -0.2\ \text{V}$ achieved $\Delta V_{\mathrm{Key0}} = 0.9747\ \text{V}$ and $\mathrm{SNM}_{\mathrm{Hold}} = 167.12\ \text{mV}$ under the authorized key, whereas $\text{Key}=1$ collapsed retention to $\mathrm{SNM} = 0.00\ \text{mV}$ ($\Delta V_{\mathrm{Key1}} = 0.0762\ \text{V}$), with \textit{Read Disturb} mitigated via \textit{Wordline Underdrive Read-Assist} ($\mathrm{SNM}_{\mathrm{Read,Assist}} = 66.23\ \text{mV}$ at $\mathrm{WL}=0.30\ \text{V}$).

Probing operating boundaries outside nominal biases within the measured LUT domain confirmed three distinct physical regimes: alongside robust nominal windows and stuck-at degeneration (\texttt{CONST\_0} at under-voltage and \texttt{CONST\_1} under asymmetric overdrive in \texttt{XOR}/\texttt{XNOR}), a spontaneous key-less \texttt{NAND} $\to$ \texttt{NOR} conversion was demonstrated at $V_{DD}=0.65\ \text{V}$ ($+0.4/-1.1\ \text{V}$), demarcating voltage margins relevant to fault-injection resilience. Across the seven-tuple DPA Pareto frontier, \texttt{(XOR, XNOR)} ($0.00\%$), \texttt{(MAJ3, XOR3)} ($0.04\%$), and \texttt{(NOR, XOR)} ($1.16\%$) stood out against the $163.91\%$ peak exposure of \texttt{(NAND, NOR)}. At the system level on a 1-bit Full Adder, replacing the 3-gate \texttt{NAND}/\texttt{NOR} carry tree with a single symmetric \texttt{(MAJ3, XOR3)} cell reduced carry-stage peak DPA asymmetry from $113.83\%$ to $0.04\%$, saved 2 transistors, increased wrong-key corruption to $87.50\%$, and lowered critical PDP by $3.01\times$, capabilities also exposed interactively through \textit{BP-FET Studio}.

With respect to methodological limitations inherited from the experimental lookup tables, the library \texttt{bp\_\allowbreak ari\_\allowbreak ngspice.lib} preserves the quasi-static formulation of \texttt{bp\_\allowbreak ari.va} without intrinsic parasitic capacitances ($C_{gs}$, $C_{gd}$, $C_{g2s}$), temperature dependence, or process variability ($\sigma_{V_{fb}}, \sigma_W$); consequently, propagation delays, switching energies, and absolute PDP values depend directly on the explicit testbench capacitances ($0.05\text{--}1.0\ \text{fF}$), while the unmeasured interval $V(g2,s) \in (-0.90,\allowbreak\ +0.20)\ \text{V}$ remained excluded to avoid ungrounded extrapolation. Future work includes embedding parasitic capacitive networks and statistical process parameters into the \texttt{ARI} subcircuit, scaling the 6T SRAM into $2 \times 2$ and $4 \times 4$ arrays with \textit{Write Margin} and sneak-path analysis, and extending the low-DPA polymorphic adder to multi-bit arithmetic datapaths and cryptographic primitives evaluated against SAT and correlation power analysis (CPA) attacks.
